\section{Context-free grammars and pushdown automata}
\label{sec:context-free}

This section treats a grammar as finite input data.  A variable or a terminal
is a natural number carrying a tag, a rule is a finite list, and a derivation
is an explicitly supplied finite certificate.  No set of all derivable words
is an object of arithmetic.  The notation $w\in L(G)$ abbreviates the
first-order assertion that a checked derivation certificate exists.  The
normalization and parsing theorems below subsequently replace this assertion
by one primitive recursive Boolean computation.  The conventional grammar,
pushdown, and pumping results provide the subject matter; the bounds and
certificate transformations given here supply their arithmetic content
\cite{Sipser2020}.

\subsection{Finite syntax and finite certificates}

\begin{definition}[Grammar and certificate conventions]
\label{cf:definition}
A context-free grammar is $G=(\Sigma,V,P,S)$, where
$\Sigma=[a]$, $V=[m]$, $m\geq1$, $S<m$, and $P$ is a finite list of
rules $A\to\beta$ with $A\in V$ and
$\beta\in(\Sigma\sqcup V)^*$.  Tags keep the two alphabets disjoint.
A sentential form is a word over this tagged alphabet.  A rewriting step
replaces one specified occurrence of $A$ by the right side of a specified
rule.  A derivation certificate lists the sentential forms, rule indices,
and replacement positions.  Its checker is a bounded scan of those lists.

A completed parse tree has a variable at its root.  A variable node carries
a rule index, and its ordered children have exactly that rule's right-side
labels.  Terminal nodes have no children.  A variable using an empty rule
also has no children; it contributes no letter to the yield.  The yield is
the left-to-right list of terminal leaves.  Trees can be coded as postorder
node tables: children have smaller indices, every nonroot node occurs in
exactly one child list, and the last node is the root.  These finite conditions
exclude sharing and disconnected components.  A bottom-up table computes
the yield of every node.  Write $\operatorname{Tree}(G,A,w,t)$ for the
Boolean assertion that $t$ is such a tree with root $A$ and yield $w$.
\end{definition}

Malformed data is assigned a fixed rejecting value.  Thus all operations in
this section are total functions on natural-number codes, even when their
mathematical specifications are stated only for well-formed input.

\begin{theorem}[Derivations and trees, with bounded translations]
\label{cf:tree-derivation}
There are primitive recursive translations in both directions between
derivations $A\Rightarrow^* w$ with terminal $w$ and completed parse trees
rooted at $A$ with yield $w$.  If every right side has length at most $r$, a
$k$-step derivation produces a tree with at most $1+kr$ nodes.  A tree with
$N$ nodes produces a leftmost derivation with at most $N$ steps, and all its
sentential forms have length at most $N$.
\end{theorem}

\begin{proof}
Begin with an unexpanded root labelled $A$.  Maintain the invariant that the
ordered frontier, counting unexpanded variables but ignoring expanded
empty-rule nodes, is the current sentential form.  The supplied replacement
position identifies one unexpanded variable on this frontier.  Attach the
specified right side as its ordered children.  This changes the frontier
by precisely the certified rewriting step and introduces at most $r$
nodes.  After $k$ steps the frontier is terminal, so every remaining frontier
node is a terminal and the tree is completed.  The node count follows by
adding the numbers of newly attached children.

In the other direction, visit variable nodes in preorder.  Expand the
leftmost unexpanded variable by the rule stored at that node, skipping
terminal nodes.  Its children occur before every later frontier variable,
so the next expansion again is leftmost.  Each variable is expanded once.
Every frontier consists of distinct nodes of the original tree, so its
length is at most $N$.  The final frontier is the yield.  Both procedures
use a finite work list; the respective loop bounds are $k$ and $N$.
Their frontier invariant is an equality between coded finite lists.
\end{proof}

\begin{lemma}[Nullable variables and bounded empty trees]
\label{cf:nullable}
The variables deriving the empty word are computed by $m$ rounds of
Boolean saturation.  The computation returns an explicit empty parse tree
for every variable it marks.
\end{lemma}

\begin{proof}
Put $E_0=\varnothing$ and obtain $E_{i+1}$ by adjoining every left side of a
rule whose right side contains no terminal and whose variables all belong
to $E_i$.  Keep previously marked variables.  Store the first rule that
marks each new variable and the already stored certificates for its children.
At a strict stage the number of marked variables increases.  An equality
stage is permanent, because applying the same deterministic update to the
same set gives the same result.  Consequently $E_m=E_{m+1}$: otherwise the
first $m+1$ updates would be strict in an $m$-element set.

The stored rules give empty trees by copying the previously stored child
trees.  Conversely, scan any completed empty tree in postorder.  A terminal
leaf would contribute a letter, so none occurs.  When its children have
been shown to lie in the fixed point $E_m$, the rule at a variable node
puts that node's label in $E_m$ as well.  In particular its root is marked.
If $r=\max(1,\max_{P}|\beta|)$, each stored certificate has height at most
$m+1$ and fewer than $(r+1)^{m+2}$ nodes.  This generous bound covers unary
rules and empty right sides.  Thus certificate production itself is
primitive recursive, without an unbounded search for a derivation.
\end{proof}

\begin{theorem}[Productivity, emptiness, and a bounded positive witness]
\label{cf:emptiness}
There is a primitive recursive decision procedure for $L(G)=\varnothing$.
If the language is nonempty, it produces a terminal word and a parse tree
for it of fewer than $(r+1)^{m+2}$ nodes.
\end{theorem}

\begin{proof}
Use the preceding saturation with one change: terminals in a right side
are allowed and require no membership test.  A rule becomes usable when
all of its variable occurrences are already marked.  Record a terminal
leaf for each terminal occurrence and the stored tree for each variable
occurrence.  The same increase argument gives a fixed point after $m$
rounds and the same bound gives a tree for each marked variable.  Conversely,
the postorder scan of any terminal-yield tree marks all its variable labels
in that fixed point.  The start variable is marked exactly when a derivation
certificate exists.  A negative answer therefore refutes every supplied
certificate; it does not assert failure of an unbounded search.
\end{proof}

\subsection{Normalization with witness-preserving transformations}

\begin{theorem}[Eliminating empty rules]
\label{cf:epsilon-removal}
One can transform $G$ primitively recursively into an equivalent grammar
whose only possible empty rule belongs to a fresh start variable $S_0$, and
in which $S_0$ occurs on no right side.  Both implications of language
equivalence are witnessed by primitive recursive tree transformations.
\end{theorem}

\begin{proof}
First introduce $S_0\to S$.  Compute the nullable variables.  For each rule
$A\to X_1\cdots X_k$, enumerate the $2^k$ masks of its positions.  A mask may
delete a position only when it contains a nullable variable.  Retain every
resulting nonempty right side.  Retain $S_0\to\eps$ exactly when $S$ was
nullable.  Record, with each new rule, one original rule and one deletion
mask that produced it; a bounded first-match search supplies this provenance.
No rule introduces $S_0$ on its right side.

Given an old tree with nonempty yield, process it bottom-up and delete each
child subtree with empty yield.  Every such child is a nullable variable,
so the surviving children are an enumerated new rule.  Nodes with positive
yield retain at least one child.  The root therefore yields a new tree with
the same word.  An old empty tree at the start is replaced by the recorded
rule $S_0\to\eps$.

Conversely, for each new rule use its provenance and restore every erased
child with the fixed nullable tree from Lemma~\ref{cf:nullable}.  Recurse
over the retained variable children.  The inserted trees have empty yield,
so the result has the same terminal yield.  For $S_0\to\eps$ use its stored
old empty tree.  This is a bounded traversal of the supplied tree, with a
grammar-dependent bound on the size inserted at every node.
\end{proof}

\begin{theorem}[Eliminating unit rules]
\label{cf:unit-removal}
Unit rules $A\to B$ can be removed primitively recursively, preserving
the language and the preceding convention on the start and empty rule.
\end{theorem}

\begin{proof}
For the current $M$ variables form the directed graph of unit rules.
Compute its reflexive transitive closure by paths of length at most $M-1$.
This bound is sufficient: on any longer path, find two equal vertices by
bounded search and erase the segment between them; at most the original
path length many erasures leave a path with no repeated vertex.  Record a
short path for each reachable pair.

For every reachable pair $(A,B)$ and every nonunit rule $B\to\gamma$, add
$A\to\gamma$.  In an old completed tree, contract each maximal chain of
unit nodes into one application of its last, nonunit rule.  Each such chain
must end because the input is a finite completed tree.  The resulting rule
is among the added rules.  In a new tree, expand each rule using the stored
short unit path followed by its old nonunit rule.  These operations preserve
yield.  The only empty rule before this operation belonged to $S_0$; since
no unit edge enters $S_0$, no other variable acquires an empty rule.
\end{proof}

\begin{theorem}[Effective Chomsky normal form]
\label{cf:cnf}
There is a primitive recursive function $G\mapsto C(G)$ with
$L(C(G))=L(G)$ such that every rule of $C(G)$ is
\[
 A\to BC,\qquad A\to a,\qquad\text{or, exceptionally, } S_0\to\eps,
\]
and $S_0$ never occurs on a right side.  The function returns tree
translations certifying both directions of this equality.
\end{theorem}

\begin{proof}
Apply the preceding two transformations.  In every right side of length
at least two, replace a terminal $a$ by a fresh variable $T_a$ with rule
$T_a\to a$.  Introduce at most one $T_a$ for each letter.  A long right side
$X_1\cdots X_k$, $k\geq3$, is replaced by
\[
 A\to X_1D_1,\quad D_1\to X_2D_2,\quad\ldots,\quad
 D_{k-2}\to X_{k-1}X_k,
\]
using variables private to that rule.  Expansion and contraction of these
chains, and insertion and removal of the terminal wrappers, are mutually
yield-preserving operations on the relevant tree pieces.  Unit and empty
rules are not reintroduced.  Private auxiliary variables cannot interact
with other rules, which proves the reverse implication as well.

Here is an explicit, deliberately loose size bound.  If the input has $m$
variables, $a$ letters, $p$ rules and maximum right-side length at most
$r\geq1$, put
\[
 M=m+1,\qquad K=(p+2)2^{r+1},\qquad J=M^2K.
\]
Empty-rule elimination has at most $K$ rules, unit elimination at most
$J$, and the final grammar has at most $M+a+Jr$ variables and at most
$a+Jr$ rules.  Every transformation is a finite enumeration within these
bounds.  For example, put $H=(r+1)^{M+2}$.  An original tree with $N$ nodes
can be transformed forward using at most $2(r+1)(N+1)$ nodes: addition of
the fresh start adds one, deletion and contraction reduce size, terminal
wrappers at most double it, and binary chains multiply it by at most
$r+1$.  In the reverse direction,
\[
             (M+1)(1+rH)N+H+1
\]
is a sufficient bound.  Contracting private chains and terminal wrappers
does not increase size; restoring unit paths multiplies it by at most
$M+1$; restoring deleted children inserts at most $r$ nullable trees of
size at most $H$ per variable node.  The final $H+1$ also covers the
start's empty alternative and removal of the fresh start.  Thus the
certificate translations have explicit primitive recursive size bounds.
\end{proof}

The theorem is a language-preserving normalization.  It does not claim to
preserve the number of derivations or the ambiguity of the original grammar.
For example, unit or empty cycles may produce arbitrarily many different
old trees for one word and disappear during normalization.

\subsection{Parsing and other finite decisions}

\begin{theorem}[CYK membership, reconstruction, and its arithmetic bound]
\label{cf:cyk}
Membership in a finite context-free grammar is uniformly primitive
recursive in the grammar and the input word.  For a normal-form grammar
with $M$ variables and $P$ rules, a word of length $n>0$ is decided using
at most $M(n+1)^2$ Boolean table cells and at most
$(P+1)M(n+1)^3$ elementary rule-and-split inspections.  On acceptance a tree
with $3n-1$ nodes is reconstructed, using exactly $2n-1$ rule applications.
\end{theorem}

\begin{proof}
Let $w=w_0\cdots w_{n-1}$.  For $0\leq i<j\leq n$ and a variable $A$, fill
$D(A,i,j)$ in increasing order of $j-i$.  At length one the value is the
Boolean disjunction of $A\to w_i$ over the finite rule list.  At greater
length it is
\[
 \bigvee_{A\to BC\in P}\ \bigvee_{i<k<j}
       \bigl(D(B,i,k)\wedge D(C,k,j)\bigr).
\]
Keep the least successful rule and split as a backpointer.  All referenced
intervals are shorter than the interval being filled.  The displayed cell
and inspection bounds follow simply by allowing all $M$ variables, all
$(n+1)^2$ interval endpoints, all $P$ rules, and all $n+1$ splits.

A successful terminal cell gives a one-letter tree.  A successful binary
cell joins the two trees supplied by its earlier cells.  Induction through
the table therefore constructs a valid tree with the advertised substring
as yield.  Conversely, take any valid tree for a nonempty word.  Below the
root no empty production is possible, because only $S_0$ has one and $S_0$
is absent on right sides.  At a binary node the two child yields are
nonempty.  Their boundary supplies one of the inspected splits, and a
postorder scan of the supplied tree shows that the corresponding cell is
true.  Its root supplies $D(S_0,0,n)$.

In a positive-yield normal-form tree, write $b$ for the number of binary
variable nodes and $t$ for the number of terminal-producing variable nodes.
Counting edges between variable nodes gives $2b=b+t-1$, hence $b=t-1$.
There is one terminal leaf for each terminal-producing variable, so $t=n$.
The tree thus has $(n-1)+n+n=3n-1$ nodes and $2n-1$ variable nodes.  Expanding
the backpointers with an explicit task stack constructs this many nodes;
it never searches for a tree.  For $n=0$, inspect the single permitted empty
rule.  For a general grammar, first apply Theorem~\ref{cf:cnf} and translate
the resulting certificate back.  A negative table value refutes every
supplied old certificate via the forward translation.
\end{proof}

\begin{theorem}[A direct bounded fixed-point parser for any CFG]
\label{cf:chart}
There is also a primitive recursive parser that works directly on the
original grammar, including its empty and unit cycles.  For an input of
length $n$ and $M$ variables, its chart has
$Q=M(n+1)(n+2)/2$ possible entries and stabilizes within $Q$ simultaneous
saturation rounds.
\end{theorem}

\begin{proof}
Use entries $(A,i,j)$ for $0\leq i\leq j\leq n$, allowing empty intervals.
Initially no entry is true.  At each round keep all true entries and add
$(A,i,j)$ if a rule $A\to X_1\cdots X_k$ has a nondecreasing cut vector
\[
                  i=t_0\leq t_1\leq\cdots\leq t_k=j
\]
whose pieces match its symbols.  A terminal occurrence requires
$t_\ell=t_{\ell-1}+1$ and the matching input letter.  A variable occurrence
$B$ requires the entry $(B,t_{\ell-1},t_\ell)$ from the previous round.
For an empty right side the condition is simply $i=j$.  Every rule and
cut vector lies in a finite enumeration bounded by the grammar and $n$.
An equality of two successive charts is permanent; every strict round
adds an entry.  Hence the chart after $Q$ rounds is fixed.

When an entry first appears, record its round, rule, and cut vector.
The child-variable entries used for it have strictly earlier ranks.
Recursively joining their stored trees, with terminal leaves at the
terminal positions, gives a genuine tree for the indicated substring.
The decreasing ranks prevent cycles in reconstruction even when the
grammar has cycles.  If $r\geq1$ bounds right-side lengths, fewer than
$(r+1)^{Q+2}$ nodes suffice for each reconstructed tree.

Conversely, scan any supplied parse tree in postorder, annotating each
node by its interval of terminal-leaf positions in the whole yield.
The children of a node give a nondecreasing cut vector, including equal
cuts for empty yields.  If all child-variable entries are in the final
chart, its fixed-point property includes the parent entry.  Thus every
node, including the root, is present.  The root entry $(S,0,n)$ therefore
is true exactly for generated words.  All annotations and the processed
node invariant are bounded finite data.  This proves an independent
membership algorithm without appealing to normal form.
\end{proof}

\begin{proposition}[Counting parses of a normal-form grammar]
\label{cf:parse-count}
For a normal-form grammar, the number of parse trees for a specified word
is a primitive recursive natural number.  In particular, zero, one, or
more than one parses of that word can be distinguished effectively.
\end{proposition}

\begin{proof}
Use the same interval order, replacing the Boolean disjunctions by sums and
the conjunctions by products.  At length one count the matching rule
indices; at greater length use
\[
 C(A,i,j)=\sum_{A\to BC\in P}\ \sum_{i<k<j}
                  C(B,i,k)C(C,k,j).
\]
A tree of the given interval has a unique root rule, a unique boundary
between its child yields, and an ordered pair of child trees.  Conversely
any such data joins to a tree.  The finite disjoint union and product count
is therefore exact.  Rules may either be deduplicated first or treated as
distinct indexed productions consistently with the tree convention.
There are only $3n-1$ nodes in a tree for $n>0$, which also directly gives
a finite enumeration bound.  Empty input is handled by counting the
start's empty rules.  This proposition concerns the normal-form grammar,
not the number of parses of an arbitrary grammar before normalization.
\end{proof}

\subsection{Closure constructions}

When a finite list of factor words is supplied with membership assertions
instead of factor trees, Theorem~\ref{cf:cyk} supplies the latter uniformly:
run its primitive recursive certificate constructor on every factor.
Its completeness theorem proves each resulting tree valid, and its
explicit size bound bounds the entire finite list of trees.  Thus the
tree-list constructions below also establish the ordinary membership
formulations of star and substitution.  They do not use an unbounded
finite-collection or choice principle for existential derivation witnesses.

\begin{theorem}[Union, concatenation, and star]
\label{cf:basic-closure}
Finite CFGs can be constructed primitively recursively for
$L(G_1)\cup L(G_2)$, $L(G_1)L(G_2)$, and $L(G_1)^*$.  Tree and word-factor
certificates are transformed in both directions.
\end{theorem}

\begin{proof}
Rename variable sets into disjoint finite intervals.  For union add a fresh
start $T$ and $T\to S_1\mid S_2$; its first rule selects exactly one of the
two old trees.  For concatenation use $T\to S_1S_2$.  Its two ordered subtrees
give a factorization $w=uv$, and any supplied pair of trees for $u,v$ joins
under this rule.

For star add $T\to\eps\mid S_1T$.  Repeatedly inspect the second child of a
nonempty-rule $T$ node.  Because the tree is finite, this process ends at
$T\to\eps$ after at most its node count many steps.  The first children give
a finite list of old-language factors, possibly empty, whose concatenation
is its yield.  Conversely, a supplied finite list of factor trees builds
this chain by recursion on the list length.  Empty factors cause no
termination problem: termination is measured by the supplied list or tree,
not by the length of its yield.
\end{proof}

\begin{theorem}[Reversal]
\label{cf:reversal}
Replacing the right side of every rule by its reversal gives a CFG for
$L(G)^R$.
\end{theorem}

\begin{proof}
Reverse the children at each variable node and perform the same operation
in every child subtree.  The resulting rule is a rule of the reversed
grammar.  A bottom-up calculation uses
$(x_1\cdots x_k)^R=x_k^R\cdots x_1^R$ to show that the yield is reversed.
Applying the operation twice restores the original tree.  It preserves
the node count and is implemented by one bounded traversal.
\end{proof}

\begin{theorem}[Arbitrary string homomorphisms]
\label{cf:homomorphism}
Let $h:\Sigma^*\to\Delta^*$ be specified by a finite table of words
$h(a)$, with empty images allowed.  Replacing each terminal occurrence
$a$ in a rule by the terminal block $h(a)$ gives a CFG for $h(L(G))$.
\end{theorem}

\begin{proof}
Retain each rule's old index as provenance.  In an old tree, replace each
terminal block in a variable node's rule by its image block, recursively
transforming variable children.  The new yield is the concatenation of the
letter images in the old yield, namely $h(w)$.

Conversely, inspect a new tree node and its stored old rule.  That old rule
specifies exactly which child positions arose from variable occurrences
and which consecutive terminal positions arose from each $h(a)$.  Replace
each latter block by one old terminal leaf $a$, even if that block has
length zero.  Recurse on the variable positions.  This gives an old tree
with yield $v$ and proves that the new yield is $h(v)$.  Ambiguity between
different old rules with the same image is resolved by the recorded rule
index, using no choice principle.  If $H=\max_a|h(a)|$, each original
terminal occurrence creates at most $H$ terminal nodes, so the tree-size
increase has an elementary bound in $H$ and the original size.
\end{proof}

\begin{theorem}[Context-free substitution]
\label{cf:substitution}
Suppose a CFG $H_a$ over $\Delta$ is supplied for every $a\in\Sigma$.
There is a CFG for all words $v_1\cdots v_k$ where
$a_1\cdots a_k\in L(G)$ and $v_i\in L(H_{a_i})$ for all $i$.
\end{theorem}

\begin{proof}
Rename all variable sets disjointly.  In the rules of $G$, replace a terminal
$a$ by the start variable of $H_a$, and include every rule of every $H_a$.
Substitute a supplied $H_a$ tree for every terminal leaf of a supplied $G$
tree to obtain a new tree with the desired concatenated yield.

For the reverse transformation, a variable from a component $H_a$ can
only have descendants from that component and terminals of $\Delta$;
none of its rules returns to a $G$ variable.  In a new tree, contract each
maximal component subtree entered from a $G$ rule to the old terminal $a$.
The contracted tree is an old $G$ tree, and the contracted pieces are
exactly the required factor trees.  The boundaries and factor order are
found by a finite scan of the tagged nodes.  This argument also handles
empty component yields and empty source words.
\end{proof}

\begin{theorem}[Intersection with a regular language]
\label{cf:regular-intersection}
Given a CFG $G$ and a DFA $D$ with $d$ states, one can construct a CFG for
$L(G)\cap L(D)$ primitively recursively.  If $G$ is in normal form with
$M$ variables, the product has at most $Md^2+1$ variables.
\end{theorem}

\begin{proof}
Use variables $[p,A,q]$, intended to generate words derived from $A$ that
take the DFA from $p$ to $q$.  For every binary rule add
\[
 [p,A,q]\to[p,B,r][r,C,q]\qquad(A\to BC),
\]
for all $p,r,q$.  For $A\to a$ add
$[p,A,\delta(p,a)]\to a$.  If $S_0\to\eps$ is present add
$[p,S_0,p]\to\eps$.  A fresh start $T$ has rules
$T\to[q_0,S_0,f]$ for each final state $f$.

Given an old tree and a DFA run on its yield, assign to every subtree its
entry and exit states in that run.  The boundary between the two child
yields gives $r$, so this decoration is a new parse tree.  Conversely,
erase the state tags in a new tree.  At every node the recorded endpoints
compose exactly as DFA runs compose under concatenation.  A bottom-up
scan thus verifies the intended interpretation of $[p,A,q]$; at the fresh
start it gives a run from $q_0$ to a final state.  Empty yield is handled
by the displayed diagonal rule.  All state assignments and rule generation
range over explicit finite tables.
\end{proof}

\begin{corollary}[Disjointness and a one-sided inclusion decision]
\label{cf:regular-decisions}
For a CFG $G$ and a DFA $D$, disjointness
$L(G)\cap L(D)=\varnothing$ and inclusion $L(G)\subseteq L(D)$ are primitive
recursively decidable.  Equality of a CFG language with a supplied finite
language is also primitive recursively decidable.
\end{corollary}

\begin{proof}
For disjointness apply the preceding construction and the productivity
test.  For inclusion replace the DFA's final-state set by its complement
and test emptiness of the resulting intersection.  A negative result
returns an actual counterexample word by Theorem~\ref{cf:emptiness}.
For a finite list $F$, construct its finite trie automaton, test
$L(G)\subseteq F$ in this way, and test each $w\in F$ by CYK for the other
inclusion.  Every membership predicate involved is Boolean, so these
finite tests justify the stated alternatives intuitionistically.  This
does not give a decision procedure for $L(D)\subseteq L(G)$ for an arbitrary
regular language, or for equality of two arbitrary CFG languages.
\end{proof}

\begin{theorem}[A finite EBNF-to-CFG compiler]
\label{cf:ebnf}
Consider finite grammar definitions whose right sides are finite expression
trees built from terminal letters, references to named variables, the empty
word, binary alternation and concatenation, optional expressions, and
star or positive repetition.  Interpret recursive references by finite
parse certificates.  There is a primitive recursive translation to a CFG
with the same terminal yields, with certificate transformations in both
directions.  If there are $m$ named variables, $p$ definitions, and $e$
expression-node occurrences, the result has $m+e$ variables and at most
$p+2e$ rules, all of right-side length at most two.
\end{theorem}

\begin{proof}
Give every expression-node occurrence $E$ a fresh private variable $X_E$.
For a definition $A:=E$, add $A\to X_E$.  Compile the expression nodes by
the following finite table, where $F,H$ denote their child occurrences:
\[
\begin{array}{c|l}
 E & \text{rules for }X_E\\ \hline
 \eps & X_E\to\eps\\
 a & X_E\to a\\
 \text{reference }B & X_E\to B\\
 F\mid H & X_E\to X_F\mid X_H\\
 FH & X_E\to X_FX_H\\
 F\text{ optional} & X_E\to\eps\mid X_F\\
 F^* & X_E\to\eps\mid X_FX_E\\
 F^+ & X_E\to X_F\mid X_FX_E
\end{array}
\]
Each occurrence contributes one variable and at most two rules, giving
the stated bounds.  Distinct tags keep expression occurrences, original
variables, and terminal letters separate even when their printed names
coincide.  Construction is a bounded scan of the finite syntax tree.

Here the source certificate has explicit constructors: a named variable
chooses one of its definitions; a terminal or empty expression has its
specified yield; a reference carries a certificate for the named variable;
alternation records its chosen branch; concatenation carries an ordered
pair; an optional expression records absence or one child; a star carries
a finite list of child certificates; and a positive repetition carries
a nonempty such list.  This is a finite tagged tree with finite list
fields.  Its local constructor conditions and yield equations are checked
in postorder, just as for CFG trees.  Cyclic references do not change the
definition: only an actually supplied finite certificate is accepted.

Translate a source certificate according to the table.  A star list
becomes the rightward chain of $X_E\to X_FX_E$ rules, ending in
$X_E\to\eps$.  For a nonempty repetition list, use that chain up to its
last member and end in $X_E\to X_F$.  The other constructors correspond
directly to the displayed rules.  The concatenation of the child yields
is unchanged at every node, so the resulting CFG tree has the same word.

Conversely, inspect a completed CFG tree.  Every private expression
variable has exactly the rules prescribed for its occurrence, so its
rule index identifies the original constructor.  For a star or positive
repetition, scan its $X_E$ rightward chain to obtain the finite list of
$X_F$ subtrees.  The chain ends within the supplied tree's node count;
positive repetition cannot terminate by an empty rule.  Recursively
translate these subtrees and the other constructor children.  Original
variable nodes recover the selected definition.  This constructs a valid
source certificate with the same yield.

Both traversals are bounded by their supplied certificate tables; every
source node or list cell emits a bounded number of new nodes.  If only
source-language factor membership assertions are given, the compiled
grammar's primitive recursive parser supplies the finite witnesses as
above.  Thus optionality and repetition introduce no infinitary proof
principle.  Ordinary BNF is the fragment using named variables, terminal
strings, alternatives, and explicit empty right sides; its CFG interpretation
is exactly the corresponding part of this construction.
\end{proof}

\subsection{Pushdown automata and inverse homomorphisms}

\begin{definition}[The exact pushdown convention]
\label{cf:pda-definition}
A nondeterministic PDA has a finite state set $Q$, a finite stack alphabet
$\Gamma$, an initial state $q_0$, one initial stack symbol $Z_0$, and a
finite transition list.  A transition
$(p,b,X)\mapsto(q,\beta)$ consumes $b\in\Sigma\cup\{\eps\}$, removes
the top symbol $X$, and replaces it by $\beta\in\Gamma^*$, written with
its new top on the left.  No transition is available on an empty stack.
Acceptance means: some supplied finite run consumes the entire input and
finishes with empty stack, in any state.  Epsilon transitions consume no
input.  A run checker examines a finite sequence of configurations and
transition indices; it is primitive recursive.
\end{definition}

\begin{theorem}[From CFGs to PDAs]
\label{cf:cfg-to-pda}
A one-state PDA in this convention can be constructed from a CFG, with
exactly the same language.  Finite acceptance certificates translate in
both directions primitively recursively.
\end{theorem}

\begin{proof}
The stack alphabet is the tagged union of variables and terminals, and the
initial stack is $S$.  For $A\to\beta$ use the epsilon transition replacing
top $A$ by $\beta$.  For each terminal $a$, use the transition consuming
$a$ and removing top $a$.  A preorder traversal of an old parse tree expands
its variable nodes and then matches its terminal leaves in order.  It uses
one transition per variable node or terminal leaf, and ends with empty
stack after reading exactly the yield.

In any accepting PDA run maintain the sentential form consisting of the
already consumed terminal prefix followed by the current stack.  An
epsilon expansion is exactly an old grammar rewriting step at the first
remaining variable position.  A terminal-matching step leaves this whole
sentential form unchanged, moving its split point one place to the right.
At the end it is the input word.  Deleting the terminal-matching steps
therefore gives a grammar derivation, which Theorem~\ref{cf:tree-derivation}
converts to a tree.  Both verifications are bounded scans of the given
tree or run.
\end{proof}

\begin{theorem}[From PDAs to CFGs, without a push-length assumption]
\label{cf:pda-to-cfg}
For every PDA in Definition~\ref{cf:pda-definition}, a finite CFG with the
same language and primitive recursive certificate translations can be
constructed.  If there are $q$ states, $g$ stack symbols, $t$ transitions,
and no transition pushes more than $r$ symbols, the grammar has
$q^2g+1$ variables and at most $q+t(q+1)^{r+1}$ rules.
\end{theorem}

\begin{proof}
Use variables $[p,X,s]$.  They describe a run segment that starts in state
$p$ with stack $X\gamma$, finishes in state $s$ with stack $\gamma$, and
does not expose $\gamma$ before its last configuration.  This last condition
is essential: the segment must remove this occurrence of $X$ and everything
placed above the untouched suffix.

For a transition $(p,b,X)\mapsto(u,Y_1\cdots Y_k)$ and every tuple
$s_1,\ldots,s_k\in Q$, add
\[
 [p,X,s_k]\to b\,[u,Y_1,s_1][s_1,Y_2,s_2]\cdots
                              [s_{k-1},Y_k,s_k].
\]
An epsilon $b$ contributes no terminal.  When $k=0$, instead add only
$[p,X,u]\to b$.  A fresh start has rules
$T\to[q_0,Z_0,s]$ for all $s\in Q$.  The finite rule count is bounded by
the displayed expression.

From a parse tree, execute its root transition, then the child run
segments in their order.  The first child removes $Y_1$ while leaving
$Y_2\cdots Y_k\gamma$ untouched; the next removes $Y_2$, and so forth.
Their terminal yields concatenate to exactly the consumed input of the
whole segment.  The start rule gives an accepting run.

Conversely, consider a supplied proper segment for $X$.  After its first
transition pushes $Y_1\cdots Y_k$, locate the first configuration exposing
$Y_2\cdots Y_k\gamma$, then the first exposing
$Y_3\cdots Y_k\gamma$, and so on.  These are bounded first-occurrence
searches within the supplied run.  They exist because each transition
removes only one old top symbol, so the stack height can fall by at most
one per step.  Before each first occurrence the protected suffix has not
been touched; the resulting pieces are proper segments for $Y_i$.  Their
boundary states supply precisely one enumerated grammar rule.  If $k=0$,
properness forces the segment to end immediately after its first step.
Recursing on the strictly shorter pieces constructs the tree.

This recursion can be implemented by a work list of subintervals of the
original run.  Each recursive node accounts for a different transition,
so there are no more nodes of this kind than run steps.  Terminals add at
most one leaf per transition.  With the start node, an accepting $N$-step
run gives a tree with at most $2N+1$ nodes.  Conversely the number of PDA
steps is at most the number of variable nodes of a grammar tree.  These
bounds make both transformations primitive recursive.
\end{proof}

\begin{corollary}[PDA membership with a bounded certificate constructor]
\label{cf:pda-membership}
PDA membership, including machines with epsilon cycles, is uniformly
primitive recursively decidable.  On acceptance a primitive recursive
function returns an accepting run; its length has a primitive recursive
bound in the machine and word codes.
\end{corollary}

\begin{proof}
Convert to a grammar, normalize it, and apply CYK.  Translate an accepted
normal-form tree back through the explicitly bounded constructions, then
through Theorem~\ref{cf:pda-to-cfg}.  The composition of their bounds is
primitive recursive.  A failed parsing table rules out every PDA run by
the forward certificate transformations.  The procedure never needs to
bound or traverse arbitrary epsilon computations of the original machine.
\end{proof}

\begin{theorem}[Inverse homomorphisms, including erasing maps]
\label{cf:inverse-homomorphism}
For any finite homomorphism table $h:\Sigma^*\to\Delta^*$ and CFG $G$ over
$\Delta$, a CFG for $h^{-1}(L(G))$ can be constructed primitively recursively.
\end{theorem}

\begin{proof}
Convert $G$ to a PDA $M$.  Construct a new PDA with a protected bottom marker
$\bot$ and a finite buffer in its control state.  The buffer values are
the empty word and all suffixes of the finitely many words $h(a)$; they
form a finite, explicitly enumerable set.  Initially an epsilon transition
places $M$'s initial symbol above $\bot$ and chooses its initial state with
empty buffer.

When the buffer is empty, reading an actual input letter $a$ replaces it
by $h(a)$ and leaves the simulated state and stack unchanged.  Implement
this by, for every possible top symbol including $\bot$, a transition
that replaces that symbol by itself.  A simulated epsilon transition of
$M$ is allowed regardless of the buffer.  A simulated $b$-reading transition
of $M$ is allowed exactly when the buffer begins with $b$; it consumes
that buffered letter and no actual input.  Simulated transitions are
available only above $\bot$.  When the simulated stack is empty and the
buffer is empty, an epsilon transition may pop $\bot$.  It is the only
way to empty the new stack.

During a run, the loaded images, after removing the buffered suffix,
are exactly the letters consumed by the simulated transitions.  This is
an equality of finite words checked after every step.  An accepting run
has read all original input and has empty buffer, so its simulated
transitions accept $h(w)$ in $M$.  Conversely, take an accepting run of
$M$ on $h(w)$.  Load the letters of $w$ in order whenever the buffer is
empty, interleaving the supplied simulated transitions.  Empty images
are loaded and immediately leave an empty buffer; they cause no problem
even after the simulated stack has emptied, because the bottom marker
remains.  After all input and buffered symbols have been processed, pop
the marker.  The new run has the old number of simulated steps, plus
$|w|$ loading steps and two boundary steps.  This proves both implications
by explicit run transformations.  Finally convert the new PDA to a CFG.
\end{proof}

\subsection{Pumping, finiteness, and constructive separations}

\begin{theorem}[The context-free pumping lemma with a PR witness function]
\label{cf:pumping}
Let a normal-form grammar have $M\geq1$ variables and put $p=2^M$.
From a parse tree for $w$ with $|w|\geq p$ one can primitively recursively
extract $u,v,x,y,z$ such that
\[
 w=uvxyz,\qquad |vy|>0,\qquad |vxy|\leq p,
 \qquad uv^i x y^i z\in L(G)\quad(i\in\N).
\]
The last assertion has a single primitive recursive tree constructor in
$i$.  For an arbitrary CFG the same result holds with $M$ taken from
its computed normal form, and with trees translated back to that CFG.
\end{theorem}

\begin{proof}
Consider the binary skeleton consisting only of variable nodes.  A
terminal-producing variable is a leaf.  A skeleton in which every
root-to-leaf path contains at most $h$ variable nodes has at most
$2^{h-1}$ leaves: at a binary root the two bounds add, and a leaf has one
leaf.  This calculation is proved by a bounded bottom-up scan of a given
tree, storing height and leaf count.  Since the tree for $w$ has $|w|$
such leaves, a longest path contains at least $M+1$ variable nodes.

Take the last $M+1$ variable nodes on a longest path.  Two have the same
label $A$, by finite pigeonhole; choose such a pair by bounded search,
with one a proper ancestor of the other.  The upper one's subtree has
variable height at most $M+1$.  Indeed the chosen path is longest also
inside each of its subtrees, or replacing its suffix would give a longer
whole path.  Its yield therefore has length at most $2^M=p$.

Let $x$ be the lower subtree's yield, $v$ and $y$ the terminal yields to
its left and right within the upper subtree, and $u,z$ the yields outside
the upper subtree.  These are consecutive intervals of $w$, so
$w=uvxyz$ and $|vxy|\leq p$.  Between the repeated variables there is a
binary node and hence a sibling subtree.  Every such sibling has nonempty
yield, because an empty rule cannot occur below the fresh start.  Thus
$|vy|>0$.

The piece between the two $A$ occurrences is a context with one hole
labelled $A$, producing $v$ before that hole and $y$ after it.  Delete
this context for $i=0$, use it once for $i=1$, and wrap $i$ copies around
the lower tree in general.  A direct concatenation calculation gives
$v^i x y^i$ as the new upper yield.  Each splice joins equal variable
labels, so it remains a valid parse tree.  A primitive recursion on $i$
copies a fixed finite context; if the original tree has $N$ nodes,
$(i+1)N$ is a sufficient size bound.  Its validity and yield equations
are quantifier-free induction invariants.  This constructs certificates
for every $i$, rather than inferring their existence from an infinite
set argument.
\end{proof}

\begin{theorem}[Deciding finiteness and returning either kind of witness]
\label{cf:finiteness}
Finiteness of a CFG language is primitive recursively decidable.  The
procedure returns either an exhaustive finite list of its words, or a
word with a pumping decomposition that yields accepted words of arbitrarily
large length.  If $p=2^M$ for the normalized grammar, it suffices to test
all words with
\[
                  p\leq |w|<2p.
\]
\end{theorem}

\begin{proof}
If this finite test finds an accepted word, use
Theorem~\ref{cf:pumping}.  Since $|vy|>0$, the constructed word for index
$i$ has length $|uxz|+i|vy|\geq i$.  Taking $i=N+1$ supplies an accepted
word longer than any supplied bound $N$.

Suppose the test finds none.  To refute any alleged accepted tree of length
$n\geq p$, transform it to normal form.  If its length is at least $2p$,
pump down once.  The length decreases by an integer between $1$ and $p$,
so it remains at least $p$.  Repeat this bounded process at most the
initial $n$ times, leaving the tree unchanged once its length is less
than $2p$.  The decreasing-length invariant forces arrival at an accepted
tree of length in $[p,2p)$, contradicting the finite test.  Thus every
accepted word has length less than $p$.  Enumerate all such words and
retain precisely those accepted by CYK; this is the required exhaustive
list.  All word enumerations have a bound supplied by the finite alphabet
and $2p$, and the pumping-down argument uses the supplied tree's length
as its loop bound.  No least member of an unbounded set is invoked.
\end{proof}

\begin{proposition}[The equal-block language]
\label{cf:equal-blocks}
Over distinct letters $a,b$, the grammar $S\to aSb\mid\eps$ generates
exactly $E_{ab}=\{a^n b^n:n\in\N\}$.  Thus, together with its regular
pumping obstruction, it witnesses the strict inclusion of regular
languages in context-free languages.
\end{proposition}

\begin{proof}
Applying the first rule $n$ times and the second rule once is a certificate
for $a^n b^n$.  Conversely, in any completed tree each nonempty rule has
exactly one variable child.  Follow this chain until its final empty rule;
its finite number $n$ of nonempty rules supplies exactly $n$ initial $a$'s
and $n$ final $b$'s.  The traversal is bounded by the given tree size.
For completeness, if a $d$-state DFA accepted precisely this language,
the first $d+1$ states while reading the prefix $a^d$ would repeat.
Deleting the nonempty loop from the accepted word $a^d b^d$ preserves
acceptance but changes the number of $a$'s, a contradiction.  The loop is
found by finite pigeonhole.  A DFA has a nonempty state set, so $d\geq1$.
\end{proof}

The balanced-parenthesis grammar $S\to\eps\mid(S)S$ gives another
important example.  Its exact balance characterization, forced first-return
split, unique parse tree, and enumeration are proved together in
Theorem~\ref{count:dyck}, avoiding separate parsing and counting assumptions.

\begin{theorem}[An explicit separating word against every CFG]
\label{cf:abc-separator}
Let $T=\{a^n b^n c^n:n\in\N\}$ over three distinct letters.  There is a
primitive recursive algorithm which, on every finite CFG $G$ over this
alphabet, returns a word where $L(G)$ and $T$ differ.  With $M$ variables
in $C(G)$, the word has length at most $3\cdot2^M$.  In particular $T$
is not context-free, in the constructive sense that every proposed
grammar is effectively refuted.
\end{theorem}

\begin{proof}
Set $p=2^M$ and run the parser on $w=a^p b^p c^p$.  If it rejects, $w$ is
the required separating word, with its negative parsing table as
certificate.  If it accepts, extract a tree and the pumping decomposition
$w=uvxyz$.  The interval $vxy$ has length at most $p$.  An interval
containing both an $a$ and a $c$ would contain all $p$ intervening $b$'s
and at least those two additional letters, so it would have length at
least $p+2$.  Therefore $v$ and $y$ together affect at most two of the
three letter types.  They remove at least one letter because $|vy|>0$.

The pumped-down word $uxz$ is still accepted.  One letter type retains
its original count $p$, while another has smaller count.  Consequently
the three counts are unequal, so $uxz\notin T$.  Its length is at most
$3p$, and its acceptance certificate is supplied by pumping.  These two
branches are selected by one Boolean computation, not by excluded
middle for an arbitrary language predicate.  The membership and
nonmembership assertions for $T$ are just finite order and count checks.
\end{proof}

\begin{theorem}[Failure of intersection and complement closure]
\label{cf:nonclosures}
There are explicit context-free languages whose intersection is not
context-free, and an explicit context-free language whose complement
relative to $\{a,b,c\}^*$ is not context-free.
\end{theorem}

\begin{proof}
Let
\[
 L_1=\{a^n b^n c^k:n,k\in\N\},\qquad
 L_2=\{a^k b^n c^n:n,k\in\N\}.
\]
For $L_1$ use start $U\to SC$, with
$S\to aSb\mid\eps$ and $C\to cC\mid\eps$.
For $L_2$ use start $U\to AR$, with
$A\to aA\mid\eps$ and $R\to bRc\mid\eps$.
The chain argument of Proposition~\ref{cf:equal-blocks} verifies their
languages, and a word belongs to both exactly when all three block
lengths agree.  Their intersection is $T$, excluded by
Theorem~\ref{cf:abc-separator}.

One can exhibit the complement counterexample directly, avoiding any
appeal to a general classical De Morgan principle.  Write $E_{bc}$ for
the equal-block grammar over $b,c$, and let
\begin{align*}
 R_0&=\{a,b,c\}^*\setminus a^*b^*c^*,\\
 K&=R_0\ \cup\ (a^+E_{ab}\cup E_{ab}b^+)c^*
       \ \cup\ a^*(b^+E_{bc}\cup E_{bc}c^+).
\end{align*}
The language $R_0$ is regular; all the other displayed pieces are
context-free by the already proved union and concatenation constructions.
They form an explicit CFG expression, and those constructions compile
it to a finite grammar.  A word outside block order belongs to $R_0$.
A word in block order is $a^i b^j c^k$; its counts are not all equal
exactly when $i>j$, $i<j$, $j>k$, or $j<k$.  These four decidable cases
are exactly the other four pieces of $K$.  Hence $K$ is precisely the
complement of $T$, and its complement is $T$.  All equivalences here
are verified using Boolean order and count comparisons.
\end{proof}

\begin{proposition}[A bounded tape recognizer and the hierarchy separation]
\label{cf:abc-noncontracting}
The language $T$ has a deterministic recognizer using only the input's
cells, fixed endmarkers, and finitely many marked variants of its letters.
The computation has a primitive recursive quadratic time bound.  Moreover,
context-free languages are a proper subclass of the languages generated
by noncontracting grammars with the isolated-start empty-word convention.
\end{proposition}

\begin{proof}
First scan to check block order.  Repeatedly mark the leftmost unmarked
$a$, then the leftmost unmarked $b$, then the leftmost unmarked $c$,
rejecting if either latter search fails.  When no unmarked $a$ remains,
accept exactly if no unmarked $b$ or $c$ remains.  Each successful round
marks one of each letter.  At most the input length $n$ many rounds are
possible, and each search and return to the left end is a bounded tape
scan.  Fewer than ten scans per round, and finitely many initial and
final scans, give the loose move bound $20(n+2)^2$.  Permanent endmarkers
and three marked variants of the letters suffice; no additional unbounded
tape is used.  Empty input is accepted immediately.

Counting marked letters proves soundness.  If the three block counts
are equal, after each successful round their unmarked remainders are
still equal, so the next indicated searches succeed until all are zero.
This proves completeness by a bounded invariant on the rounds.

For the language-class inclusion, the normal form from
Theorem~\ref{cf:cnf} is itself noncontracting apart from its isolated
start's empty rule.  The explicit noncontracting grammar for $T$ is
proved in Proposition~\ref{eff:abc}.  Theorem~\ref{cf:abc-separator}
excludes every CFG for that same language, giving strictness.
\end{proof}

\subsection{What the arithmetic and Lean claims mean here}
\label{cf:audit}

Every induction in this section can be organized around one of four
finite records: a processed prefix of a node table, a processed prefix
of a supplied transition trace, a bounded grammar-enumeration table, or
an already filled portion of a parsing table.  Each record's validity,
list equality, label agreement, size bound, and Boolean cell value are
primitive recursive predicates.  A conjunction over all already processed
indices is itself a bounded Boolean fold.  Quantifier-free induction on
that fold proves the invariant.  Apparent tree induction in the prose is
implemented by a postorder node loop; recursive certificate construction
is implemented by its explicitly bounded work list.  The pumping index
uses ordinary primitive recursion with a fixed finite context.

In particular, the saturation arguments do not use induction on the
unbounded statement ``$A$ has some derivation.''  The positive direction
checks a stored finite witness; the negative direction transforms or
scans an arbitrary witness supplied as a parameter.  The same distinction
justifies the negative answers of CYK and the emptiness and finiteness
procedures.  All infinite-language conclusions are expressed by a uniform
word-and-certificate constructor for each requested numerical bound.

The accompanying \texttt{lean/ContextFreeOutline.lean.txt} specifies the
finite records, Boolean predicates, algorithms, theorem interfaces, and
the order of the proof obligations.  It is an outline, not a claimed
compiled formalization.  Existing Lean grammar definitions can serve as
an external semantic comparison, but using a theorem about unrestricted
Lean predicates alone would not establish its derivability in \ipra.
The object-level route is to implement the finite functions in the
repository's primitive-recursive coding layer and prove the Boolean
invariants just described.
